
%%%%%%%%%%%%%%%%%%%%%%%%%%%%%%%%%%%%%%%%%%%%%%%%%%%%%
% Index sets

To succinctly state our theorems for double sums, it will be convenient to
introduce a more general index notation.

\begin{defn}\deflabel{index_sets}
    %
    Let $\indexset$ denote a set of sets of data indices, where each member $s \in
    \indexset$ has the same cardinality, and for $s \in \indexset$, let $\xvec_s$
    denote the set of corresponding datapoints.  We write $\phi(\xvec_s)$ for a
    function that depends on $|s|$ datapoints.  For example, if $\indexset = \{1,
    \ldots, N \} =: [N]$, then $\xvec_s = \x_s$, and $\means \phi(\xvec_s) = \meann
    \phi(\x_n)$.  Similarly, if $\indexset = [N] \times [N]$, then $\xvec_s =
    (\x_{s_1}, \x_{s_2})$, and $\means \phi(\xvec_s) = \frac{1}{N^2} \sum_{n=1}^N
    \sum_{m=1}^N \phi(\x_n, \x_m)$.
    %
\end{defn}

We now state a generalization of \defref{bclt_okay} that applies to 
generic index sets as defined in \defref{index_sets}.

\begin{defn}\deflabel{bclt_okay_index}
%
For an index set $\indexset$, a data-dependent function
$\phi(\theta, \xvec_s)$ is $K$-th order BCLT-okay if,
for each $k = 1, \ldots, K$,
%
\begin{enumerate}
%
\item\itemlabel{bclt_okay_as_index}
    $\theta \mapsto \phigrad{k}(\theta, \xvec_s)$ is BCLT-okay
    $\fdist$-almost surely
%
\item\itemlabel{prior_op1_index}
    $\means \expect{\prior(\theta)}{\norm{\phi(\theta, \xvec_s)}^2_2} = O_p(1)$.
%
\item\itemlabel{grad_op1_index}
For some $\delta > 0$,
$\sup_{\theta \in \thetaball{\delta}}
    \means \norm{\phigrad{k}(\theta, \xvec_s)}^2_2 = O_p(1)$.
%
\end{enumerate}
%
\end{defn}


Note that a function is BCLT--okay according to \defref{bclt_okay} precisely if
it is BCLT--okay according to \defref{bclt_okay_index} with index set $\indexset
= [N]$.  Although $\indexset = [N]$ and $\indexset = [N]\times[N]$ will be the
only examples of $\indexset$ we will use in the present paper, we believe it is
clearer to keep the notation general.   

In particular, \thmref{bayes_clt_main} is stated in terms of $\indexset = [N]$,
but we will prove that it holds for generic index sets of the form
\defref{bclt_okay_index}.  Formally, we will prove \thmref{bayes_clt_main} under
the following \assuref{core_bclt_assu}, noting that \assuref{core_bclt_assu} is
assumed by \thmref{bayes_clt_main} for $\indexset = [N]$.

\begin{assu}\assulabel{core_bclt_assu}
%
Assume that $\phi(\theta, \xvec_s)$ is third-order BCLT-okay
according to \defref{bclt_okay_index}, and that \assuref{bayes_clt} holds.
%
\end{assu}
